\organization{Meta}{Menlo Park, CA}
\experience{Production Engineer}{2023-12-4}{}{
    Built capacity management and model serving infrastructure for Meta's AI Inference Platform.
    \begin{itemize}
        \item Designed and built automated pre-scaling for large traffic experiments, estimating capacity needs and scaling models ahead of traffic shifts. Served 300+ quick experiments for Instagram and \href{https://en.wikipedia.org/wiki/Meta_Superintelligence_Labs}{MSL}. Extended to support pre-scaling in continuous capacity planning mode.
        \item Built launch capacity estimation and observability tooling. Implemented buffer calculation, efficiency metrics, a debugging UI, and prepare-step estimation to surface deficits before launch. Built capacity estimation for distributed inference models. Unblocked 70+ launches as oncall and drove follow-ups. Drove launch self-serve rate from 30\textendash 40\% to 60\textendash 70\%.
        \item Extended the capacity solver to support server groups and host profiles, enabling gang scheduling, distributed inference paradigms, and trusted execution environments.
        \item Migrated all inference platform configs to a config-as-code system. Co-designed and built automated config generation tooling to reduce config boilerplate.
        \item Designed and built tooling to bridge physical capacity and virtual quota. Built a capacity-to-quota syncer and shared pool accounting tools to track resource ownership and automatically allocate quota from capacity transfers, eliminating manual quota assignment.
    \end{itemize}
}

\experience{Production Engineer Intern}{2022-5-23}{2022-8-12}{
    Developed observability and monitoring tools for the AI model storage team.
    \begin{itemize}
        \item Developed a structured logging library in C++ and Python. Migrated checkpointing code from a singleton logger, decoupling logging from business logic and improving log correctness.
    \end{itemize}
}
